\documentstyle[%


righttag%


,11pt%


,amssymb%


,russian%               remove in Eng.


,izvru%                 Set izven% in Eng.


]{amsart}





%


\newcommand{\ctg}{\operatorname{ctg}}


\newcommand{\re}{\operatorname{Re}}


\newcommand{\im}{\operatorname{Im}}


\newcommand{\tg}{\operatorname{tg}}


\renewcommand{\baselinestretch}{0.96}


\newcommand{\tts}[1]{}





%\hoffset=-15mm%                  For Eng.


\setcounter{page}{3}


\god{2001}


\nomer{10~(473)} %                        Remove in Eng.


%\nomer{Vol.\,45, No.\,10}%               Set in Eng.


%\str{\pageref{begin}--\pageref{end}}%  Set in Eng.


\udk{517.544}





\author{\it Н.Р.\,Абубакиров, Р.Б.\,Салимов, П.Л.\,Шабалин}


% NB:   ^^ remove in Eng.


\title{Внешняя обратная краевая задача


при комбинировании двух параметров из


декартовых координат\\ и полярного угла}





\date{}


%\pagestyle{myheadings}%         Set in Eng.


\pagestyle{plain} %              Remove in Eng.


\begin{document}


%\thispagestyle{plain}%          Set in Eng.


\maketitle


\label{begin}








\section{Внешняя обратная краевая задача для параметров $x$, $y$}





Требуется найти в плоскости комплексного


переменного $z=x+iy$ контур $L_{z}=L_{z}^{1}\cup L_{z}^{2}$ и 


функцию $w(z)$, аналитическую в области


$D_{z}$, $\infty\in D_z$, $\partial D_z=L_z$,


по краевому условию 





$w=\varphi_{1}(x)+i\psi_{1}(x)$, $0\leq x\leq a$, на одной части


$L_{z}^{1}$,





$w=\varphi_{2}(x)+i\psi_{2}(x)$, $0\leq x\leq a$, на остальной части


$L_{z}^{1}$,





$w=\varphi(y)+i\psi(y)$, $-b\leq y\leq b$, на $L_{z}^{2}$.





Будем считать, что заданные однозначные функции 


определяют в плоскости $w$ замкнутый жордановый


контур $L_{w}$, являющийся границей односвязной конечной области $D_w$, причем 


$D_w=w(D_z)$, $\partial D_w=L_w=L_w^1\cup L_w^2$.


Предполагаем также, что начало координат 


не принадлежит $D_z$ и $w_0=w(\infty)$


фиксируем заранее \cite{TN}. 


Для определенности будем считать, что $O\in L^1_z$.





Отобразим конформно область


$E=\{|\zeta|<1\}$ в плоскости $\zeta =\rho e^{i\gamma}$ на $D_{w}$ 


функцией $w=\omega (\zeta)$ так,


чтобы $\omega (0) = w_{0}=w(\infty)$,


обозначим через $e^{i\gamma_{A}}$, $e^{i\gamma_{B}}$ 


точки окружности $\partial E$, 


соответствующие точкам $\varphi(-b)+i\psi(-b)$,


$\varphi(b)+i\psi(b)$ контура $L_{w}$, причем 


$0 <\gamma_{A} <\pi \gamma_{B}=2\pi -


\gamma_{A}$, а через $e^{i\gamma_{C}}$, $\gamma_{A} <


\gamma_{C} <\gamma_{B}$, --- 


точку, отвечающую $\varphi_{2}(0)+i\psi_{2}(0)$.





Соотношения $w=w(z)$, $w=\omega (\zeta)$ определяют функцию


$z=z(\zeta)$, отображающую область $E$ на $D_{z}$, причем


$z(0)=\infty$, и $\zeta=0$ есть простой полюс функции $z(\zeta)$.





Обозначим


\begin{equation} \label{2.0}


z(e^{i\gamma})=x(\gamma)+iy(\gamma),


\end{equation}


и из соответствия граничных значений функций $w=w(z)$, $w=\omega (\zeta)$


получим


\begin{alignat}{2}


\varphi_{2}(x)+i\psi_{2}(x)&=\omega (e^{i\gamma}),&\quad&\gamma_{A}


\leq\gamma<\gamma_{C},\notag \\


%


\label{2.1}


\varphi_{1}(x)+i\psi_{1}(x)&=\omega (e^{i\gamma}),&\quad&\gamma_{C}


\leq\gamma\leq\gamma_{B},\\


%


\varphi(y)+i\psi(y)&=\omega (e^{i\gamma}),&\quad &0\leq\gamma\leq


\gamma_{A},\ \,


\gamma_{B}\leq\gamma\leq 2\pi.\notag


\end{alignat}





Из первых двух равенств находим зависимость $x=x(\gamma)$,


$\gamma_{A}\leq\gamma\leq\gamma_{B}$, из последнего --- зависимость


$y=y(\gamma)$, $0\leq\gamma\leq\gamma_{A},\,


\gamma_{B}\leq\gamma\leq 2\pi$. Следовательно, получаем смешанную


краевую задачу для нахождения функции $z(\zeta)$:


\begin{equation} \label{2.3}


\begin{split}


\re z(e^{i\gamma})&=


x(\gamma),\quad\gamma_{A}\leq\gamma\leq\gamma_{B},\\


\im z(e^{i\gamma})&=


y(\gamma),\quad 0\leq\gamma\leq\gamma_{A},\ \,\gamma_{B}\leq\gamma


\leq2\pi.


\end{split}


\end{equation}


Будем считать, что функции $x(\gamma)$, $y(\gamma)$ удовлетворяют


условию Г\"eльдера.





Для решения этой задачи мы не можем непосредственно воспользоваться


известными результатами (\cite{Mu}, с.\,308--313), т.\,к. здесь искомая


функция имеет простой полюс в точке $\zeta=0$. Поэтому будем 


пользоваться методами, разработанными в статьях \cite{AS}, \cite{SSh}. 


Краевую задачу (3) представим в виде одного 


краевого условия


\begin{gather}\label{2.4}


\re [e^{-i\nu(\gamma)}z(e^{i\gamma})]=c(\gamma), \\


%


% \label{2.5}


\nu(\gamma)=\begin{cases}


-\pi/2, & \gamma\in [0,\gamma_{a}];\\


-\pi, & \gamma\in (\gamma_{A},\gamma_{B});\\


-3\pi/2, & \gamma\in [\gamma_{B},2\pi],


\end{cases}


\quad


c(\gamma)=\begin{cases}


-y(\gamma), & \gamma\in[0,\gamma_{A}];\\


-x(\gamma), & \gamma\in (\gamma_{A},\gamma_{B});\\


y(\gamma), & \gamma\in [\gamma_{B},2\pi].\notag


\end{cases}


\end{gather}





Теперь рассмотрим однородную задачу, соответствующую


(\ref{2.4}), и построим ее решение $F_{0}(\zeta)$, обращающееся


в нуль только в точках $e^{i\gamma_{A}},\;e^{i\gamma_{B}}\;$


и имеющее простой полюс в точке $\zeta=1$. 


При этом образом единичного круга в плоскости


$F_0$ будет верхняя полуплоскость без конечного 


разреза, лежащего на мнимой оси. Функция, 


осуществляющая это отображение, имеет вид


$$


F_{0}(\zeta)=e^{iq}\frac{(\zeta-e^{i\gamma_{A}})^{1/2}(\zeta-


e^{i\gamma_{B}})^{1/2}}{\zeta -1}.


$$


 


Произведение


\begin{equation}\label{4.2}


-ie^{-i\nu(\gamma)}F_{0}(e^{i\gamma})=|F_{0}(e^{i\gamma})|=


\frac{|e^{i\gamma}-e^{i\gamma_{A}}|^{1/2}|e^{i\gamma}-


e^{i\gamma_{B}}|^{1/2}}{|e^{i\gamma}-1|}


\end{equation}


является действительной величиной, поэтому краевое условие ({\ref{2.4})


можно записать в виде


$$


\re\frac{z(e^{i\gamma})}{-iF_{0}(e^{i\gamma})}=


\frac{c(\gamma)}{|F_{0}(e^{i\gamma})|}.


$$





Учитывая, что функция $z(\zeta)$ имеет простой полюс в точке


$\zeta =0$, находим (\cite{Ga}, с.269)


\begin{equation}\label{4.4}


\frac{z(\zeta)}{-iF_{0}(\zeta)}=\frac{1}{2\pi}\int_{0}^{2\pi}


\frac{c(\sigma)}{|F_{0}(e^{i\sigma})|}\frac{e^{i\sigma}+


\zeta}{e^{i\sigma}-\zeta}d\sigma + iB_{0} + C\zeta -


\frac{\overline{C}}{\zeta},


\end{equation}


где $C=A + iB$, $A$, $B$, $B_{0}$ -- произвольные действительные постоянные.


После перехода к пределу при $\zeta\to e^{i\gamma}$ формула примет


вид (см., напр., \cite{Ga}, с.\,59)


\begin{equation}\label{4.5}


z(e^{i\gamma})=-iF_{0}(e^{i\gamma})\bigg\{\frac{c(\gamma)}


{|F_{0}(e^{i\gamma})|}-\frac{i}{2\pi}\int_{0}^{2\pi}


\frac{c(\sigma)}{|F_{0}(e^{i\sigma})|}\ctg \frac{\sigma-\gamma}{2}


d\sigma +


i(B_{0}+2A\sin\gamma+2B\cos\gamma) \bigg\}.


\end{equation}





С учетом (\ref{4.2}) заключаем, что в формуле (\ref{4.5})


выражение в фигурных


скобках должно обращаться в нуль при $\gamma =0$, т.\,е.


должно выполняться соотношение


$$


-\frac{1}{2\pi}\int_{0}^{2\pi}


\frac{c(\sigma)}{|F_{0}(e^{i\sigma})|}\ctg \frac{\sigma}{2}d\sigma


+B_{0}+2B=0.


$$





Этими рассуждениями доказана 





\begin{thm}


Решение внешней обратной краевой задачи $($ОКЗ$)$


по параметрам $x, y$ дается формулой $(\ref{4.4})$ и 


зависит от двух произвольных постоянных $A$ и $B$.


\end{thm}





\section{Внешняя обратная краевая задача для параметров $\theta$, $y$}





$1^{\circ}$. Рассмотрим задачу,


отличающуюся от прежней только тем,


что на части $L_{z}^{1}$ контура $L_{z}$ значения аналитической


функции $w(z)$ заданы в виде


\begin{equation}\label{7.1}


w=\varphi_{*}(\theta) +i\psi_{*}(\theta),\ \;\alpha\leq\theta\leq 2\pi -


\beta,


\end{equation}


где $\theta =\arg z$ --- ветвь, непрерывная в плоскости, разрезанной по 


полуоси $x>0$, числа $\alpha$, $\beta$ удовлетворяют условиям


$0 <\alpha <\pi/2$, $0 <\beta <\pi/2$.





Будем считать, что эти значения и значения функции


$w(z)$, заданные на $L_{z}^{2} $ в виде


$$


w=\varphi (y) +i\psi (y),\ \;-b\leq y\leq b,


$$


в плоскости $w$ определяют замкнутый жордановый контур Ляпунова $L_{w}$,


состоящий из $L_{w}^{1}\;$ и $L_{w}^{2}$, причем


\begin{gather*}


\varphi_{*}(\alpha) +i\psi_{*}(\alpha)=\varphi (b)+i\psi(b), \\


\varphi_{*}(2\pi -\beta) +i\psi_{*}(2\pi -\beta)=\varphi (-b)+


i\psi(-b).


\end{gather*}





Для определенности предполагаем, что в рассматриваемой задаче начало 


координат не принадлежит $D_z$ и величина $w_0=w(\infty)$ задана.


Функции $\varphi (y)$, $\psi(y)$, $\varphi_* (\theta)$


и $\psi_*(\theta)$ будем считать 


дифференцируемыми, причем их производные не обращаются в нуль. 


Круг $E$


в плоскости $\zeta$ отобразим конформно


на область $D_{w}$, внутреннюю к $L_{w}$, функцией


$w=\omega (\zeta)$ при том же соответствии точек, что и выше.


Для граничных значений функции $z(\zeta)=w^{-1}\circ\omega(\zeta)$


будем использовать обозначения (\ref{2.0}).





Соотношения (\ref{2.1}) теперь заменяются следующими:


\begin{alignat*}{2}


\varphi_{*}(\theta)+i\psi_{*}(\theta)&=\omega(e^{i\gamma}),&\quad


&\gamma_{A} <\gamma <\gamma_{B},\\


\varphi(y) +i\psi(y)&=\omega(e^{i\gamma}),&\quad


&0\leq\gamma\leq\gamma_{A},\ \,


\gamma_{B}\leq\gamma\leq 2\pi.


\end{alignat*}





Из первого равенства находим зависимость $\theta =\theta(\gamma)$,


$\gamma_{A} <\gamma <\gamma_{B}$, причем


$\theta (\gamma_{B})=\alpha$, $\theta (\gamma_{A})=2\pi -\beta$,


из второго --- зависимость


$y(\gamma)$, $0\leq\gamma\leq\gamma_{A}$, $\gamma_{B}\leq\gamma\leq 2\pi$.


Теперь приходим к следующей краевой задаче Гильберта


для нахождения функции $z(\zeta)$:


\begin{equation}\label{8.2}


\begin{split}


\im[e^{-i\theta(\gamma)}z(e^{i\gamma})]&=0,\qquad \gamma_{A} <\gamma


<\gamma_{B},\\


\im z(e^{i\gamma})&=y(\gamma),\quad 0\leq\gamma\leq\gamma_{A},\ \,


\gamma_{B}\leq\gamma\leq 2\pi.


\end{split}


\end{equation}





Вначале найдем решение $F_{0}(\zeta)$ однородной задачи, соответствующей


(\ref{8.2}), обращающееся


в нуль только в точках $e^{i\gamma_{A}}$, $e^{i\gamma_{B}}$ и имеющее


простой полюс в точке $\zeta = 0$.





Краевое условие для $F_{0}(\zeta)$ запишем в виде


\begin{equation}\label{8.3}


\re[e^{-i\nu(\gamma)}F_{0}(e^{i\gamma})]=0,


\end{equation}


где


\begin{equation}\label{8.4}


\nu (\gamma)=\begin{cases}


-\pi/2, & \gamma\in [0,\gamma_{A}];\\


-5\pi/2 +\theta(\gamma), & \gamma\in (\gamma_{A},\gamma_{B});\\


-5\pi/2, & \gamma\in [\gamma_{B},2\pi].


\end{cases}


\end{equation}


Отсюда видно, что можно принять


\begin{equation}\label{9.1}


\arg F_{0}(e^{i\gamma})=\nu(\gamma)+\pi/2.


\end{equation}





Известные граничные значения действительной части аналитической в $E$ функции


\begin{equation}\label{9.2}


K(\zeta)=-i\ln\{F_{0}(\zeta)\zeta/[(\zeta - e^{i\gamma_{A}})^{\beta/\pi}


(\zeta - e^{i\gamma_{B}})^{\alpha/\pi}]\}


\end{equation}


обозначим через $q(\gamma)$. По формуле Шварца (\cite{Ga}, c.\,58) находим


\begin{equation}\label{9.5}


K(\zeta)=\frac{1}{2\pi}\int_{0}^{2\pi}q(\gamma)\frac{e^


{i\gamma}+\zeta}{e^{i\gamma}-\zeta}d\gamma,


\end{equation}


затем на основании (\ref{9.2}) определяем


\begin{equation}\label{9.10}


F_{0}(\zeta)=\zeta^{-1}e^{iK(\zeta)}(\zeta -


e^{i\gamma_{A}})^{\beta/\pi}(\zeta - e^{i\gamma_{B}})^{\alpha/\pi}.


\end{equation}





Граничные значения $F_0(t)$ функции $F_0(\zeta)$ удовлетворяют 


краевому условию (\ref{8.3}), поэтому


$-ie^{-i\nu(\gamma)\;}F_{0}(e^{i\gamma})\;$ является действительной


величиной, причем


\begin{equation}\label{10.4}


-ie^{-i\nu(\gamma)}F_{0}(e^{i\gamma})=|F_{0}(e^{i\gamma})|.


\end{equation}


С учетом (\ref{8.4}) по аналогии с (\ref{8.3}) краевое условие


(\ref{8.2}) запишем в виде


$$


%\begin{equation}\label{10.31}


\re [e^{-i\nu(\gamma)}z(e^{i\gamma})]=c(\gamma),


%\end{equation}


$$


где


$$


%\begin{equation}\label{10.41}


c(\gamma)=\begin{cases}


-y(\gamma), & \gamma\in [0,\gamma_{A}];\\


0, & \gamma\in(\gamma_{A},\gamma_{B});\\


-y(\gamma), & \gamma\in [\gamma_{B},2\pi].


\end{cases}


%\end{equation}


$$





Это условие, имея в виду (\ref{10.4}), представим следующим образом:


$$


%\begin{equation}\label{10.5}


\re \bigg[\frac{z(e^{i\gamma})}{-iF_{0}(e^{i\gamma})}\bigg]=


\frac{c(\gamma)}{|F_{0}(e^{i\gamma})|},


%\end{equation}


$$


затем с использованием формулы Шварца получим


\begin{equation}\label{10.6}


z(\zeta)=-iF_{0}(\zeta)\bigg[\frac{1}{2\pi}\int_{0}^{2\pi}


\frac{c(\gamma)}{|F_{0}(e^{i\gamma})|}\frac{e^{i\gamma}+\zeta}


{e^{i\gamma}-\zeta}d\gamma +iB_{0}\bigg],


\end{equation}


где $B_{0}$ --- произвольная действительная постоянная.


Отсюда 


\begin{gather}\label{11.2}


z(e^{i\gamma})=|F_{0}(e^{i\gamma})|e^{i\theta(\gamma)}\bigg[


-\frac{1}{2\pi}\int_{0}^{2\pi}\frac{c(\sigma)}{F_{0}


(e^{i\sigma})}\ctg \frac{\sigma -\gamma}{2}d\sigma + B_{0}\bigg],


\ \;\gamma\in (\gamma_{A},\gamma_{B}), \\


%


\label{11.3}


\re z(e^{i\gamma})=x(\gamma)=|F_{0}(e^{i\gamma})|\bigg[


\frac{1}{2\pi}\int_{0}^{2\pi}\frac{-c(\sigma)}{F_{0}


(e^{i\sigma})}\ctg \frac{\sigma -\gamma}{2}d\sigma + B_{0}\bigg],


\ \;\gamma\in [0,\gamma_{A}),\ \;\gamma\in(\gamma_{B},2\pi].


\end{gather}





Из формулы (\ref{11.2}) ясно, что постоянная $B_{0}$ должна


удовлетворять условию


\begin{equation}\label{11.4}


-\frac{1}{2\pi}\int_{0}^{2\pi}\frac{c(\sigma)}{F_{0}


(e^{i\sigma})}\ctg \frac{\sigma -\gamma}{2}d\sigma + B_{0} > 0,


\ \;\gamma\in (\gamma_{A},\gamma_{B}).


\end{equation}





С учетом равенства (\ref{9.10}) и результатов статьи \cite{Sa}


нетрудно проверить, что интеграл в последней формуле стремится к


$+\infty$ при $\gamma\to\gamma_{A}$ 


и при $\gamma\to\gamma_{B}$ для $\gamma\in (\gamma_{A},\gamma_{B})$.


Следовательно, неравенство (\ref{11.4})


может быть удовлетворено за счет выбора постоянной $B_{0}$.


Результаты из \cite{Sa} и (\ref{11.3}) показывают, что интеграл в


формуле (\ref{11.3}) стремится к $+\infty$ при $\gamma\to\gamma_{A}-0$,


$\gamma\to\gamma_{B}+0$, т.\,е. постоянную $B_{0}\;$ можно выбрать так, чтобы


\begin{equation}\label{11.6}


-\frac{1}{2\pi}\int_{0}^{2\pi}\frac{c(\sigma)}{F_{0}


(e^{i\sigma})}\ctg \frac{\sigma -\gamma}{2}d\sigma + B_{0} > 0,


\ \;\gamma\in [0,\gamma_{A}),\ \;\gamma\in(\gamma_{B},2\pi].


\end{equation}


Тогда согласно (\ref{11.3})


\begin{equation}\label{12.1}


x(\gamma) > 0,\quad\gamma\in [0,\gamma_{A}),\ \;\gamma\in (\gamma_{B},


2\pi].


\end{equation}


 


Отметим, что каждый из контуров $L_{z}^{1}$, $L_{z}^{2}$ не имеет


точек самопересечения. Эти контуры не будут иметь общих точек


кроме концов, и область $D_{z}$ будет однолистной,


если при выполнении условия


(\ref{12.1}) будет иметь место соотношение


\begin{equation}\label{12.2}


-\tg\beta\leq\frac{y(\gamma)}{x(\gamma)}\leq\tg\alpha,


\quad\gamma\in [0,\gamma_{A}),\ \;\gamma\in (\gamma_{B},2\pi].


\end{equation}


В самом деле, в этом случае для точек $L_{z}^{2}$ имеем


$z(e^{i\gamma})=x(\gamma)+iy(\gamma)$, $y(\gamma)/x(\gamma)=


\tg\theta(\gamma)$, $\theta(\gamma)=\arg z(e^{i\gamma})$ и


$$


-\beta\leq\theta(\gamma)\leq\alpha.


$$





Используя результаты из \cite{Sa}, покажем, что для $\gamma\in


(\gamma_{B},2\pi)$


функцию (\ref{11.3}) можно представить в виде


$$


x(\gamma)=b\ctg \alpha +|e^{i\gamma}-e^{i\gamma_{B}}|^{\alpha/\pi}


(x_{*}(\gamma)+B_{0}),


$$


где $x_{*}(\gamma)$ --- функция, 


удовлетворяющая условию Г\"eльдера при $\gamma\in 


[\gamma_{B},2\pi]$. Поэтому, выбрав $B_{0}$ так, чтобы


$$


\min_{\gamma_{B}\leq\gamma\leq2\pi}x_{*}(\gamma)+B_{0} > 0,


$$


получим


$$


x(\gamma) > b\ctg \alpha,\quad\gamma\in (\gamma_{B},2\pi).


$$


С другой стороны, $y(\gamma) < b$, $\gamma\in (\gamma_{B},2\pi)$,


следовательно,


$$


\frac{y(\gamma)}{x(\gamma)} <\tg\alpha,\ \;\gamma\in (\gamma_{B},2\pi).


$$


Аналогично, замечая, что $-y(\gamma) < b$, $\gamma\in (0,


\gamma_{A})$, выбором постоянной $B_{0}$


добьемся выполнения неравенства


$$


x(\gamma) > b\ctg\beta,\quad\gamma\in (0,\gamma_{A}),


$$


тогда


$$


-\tg\beta <\frac{y(\gamma)}{x(\gamma)},\quad\gamma\in (0,\gamma_{A}).


$$





Таким образом, доказана 





\begin{thm}


Решение внешней ОКЗ по параметрам $\theta$, $y$,


$\alpha\le\theta\le2\pi-\beta$,


дается формулой $(\ref{10.6})$ и зависит от одной 


вещественной постоянной $B_0$. Если $B_0$


удовлетворяет условиям $(\ref{11.4})$, $(\ref{11.6})$,


$(\ref{12.2})$, то искомая область 


$D_z$ является однолистной.


\end{thm}





$2^{\circ}$. Теперь рассмотрим предыдущую задачу в случае,


когда параметр $\theta$


в формуле (\ref{7.1}) изменяется в интервале $-\beta\leq\theta\leq


\alpha$, где $\alpha$, $\beta$ обозначают


то же, что и выше. В этом случае считаем,


что начало координат принадлежит $D_z$.





В отличие от предыдущего здесь имеем


$$


\varphi_{*}(-\beta)+i\psi_{*}(-\beta)=\varphi(-\beta)+i\psi(-\beta) ,\ 


\theta(\gamma_{A})=-\beta.


$$


В условии (\ref{8.3})


$$


%\begin{equation}\label{13.3}


\nu(\gamma)=\begin{cases}


-\pi/2, & \gamma\in[0,\gamma_{A}];\\


-\pi/2 +\theta(\gamma), & \gamma\in (\gamma_{A},\gamma_{B});\\


-\pi/2, & \gamma\in [\gamma_{B},2\pi].


\end{cases}


%\end{equation}


$$


В силу (\ref{9.1}) граничные значения действительной части


аналитической в круге $Е$ функции


$$


K(\zeta)=-i\ln\{F_{0}(\zeta)/[(\zeta-e^{i\gamma_{A}})^{\beta/\pi}


(\zeta-e^{i\gamma_{B}})^{\alpha/\pi}]\}


$$


известны, поэтому по формуле (\ref{9.5})


находим $K(\zeta)$


и после этого


\begin{equation}\label{13.6}


F_{0}(\zeta)=e^{iK(\zeta)}(\zeta-e^{i\gamma_{A}})^{\beta/\pi}


(\zeta-e^{i\gamma_{B}})^{\alpha/\pi}.


\end{equation}





Следовательно, искомое решение определяется формулой (\ref{4.5}),


поэтому


\begin{gather}\label{14.1}


z(e^{i\gamma})=x(\gamma)+iy(\gamma)=|F_{0}(e^{i\gamma})|e^{i\theta


(\gamma)}\bigg\lbrace\frac{1}{2\pi}\int_{0}^{2\pi}\frac


{-c(\sigma)}{|F_{0}(e^{i\sigma})|}\ctg \frac{\sigma -\gamma}{2}d\sigma+ \\


%


+B_{0}+2A\sin\gamma +2B\cos\gamma\bigg\rbrace,


\quad\gamma\in (\gamma_{A},\gamma_{B}),\notag \\


%


\label{14.2}


\re z(e^{i\gamma})=x(\gamma)=|F_{0}(e^{i\gamma})|


\bigg\lbrace\frac{1}{2\pi}\int_{0}^{2\pi}\frac


{-c(\sigma)}{|F_{0}(e^{i\sigma})|}\ctg \frac{\sigma -\gamma}{2}d\sigma+\\


%


+B_{0}+2A\sin\gamma +2B\cos\gamma\bigg\rbrace,


\quad\gamma\in [0,\gamma_{A}),\ \;\gamma\in(\gamma_{B},2\pi],\notag


\end{gather}


при этом необходимо выполнение неравенства 


\begin{equation}\label{14.3}


\frac{1}{2\pi}\int_{0}^{2\pi}\frac


{-c(\sigma)}{|F_{0}(e^{i\sigma})|}\ctg \frac{\sigma -\gamma}{2}d\sigma


+B_{0}+2A\sin\gamma +2B\cos\gamma > 0,\quad\gamma\in (\gamma_{A},


\gamma_{B}).


\end{equation}





Для обеспечения однолистности области $D_{z}$ потребуем, чтобы


линии $L_{z}^{1}$, $L_{z}^{2}$ были расположены по разные стороны от


прямой, проходящей через концы этих линий, с уравнением


$$


y=-b+\frac{2}{\ctg \alpha -\ctg \beta}(x-b\ctg \beta).


$$


Поэтому будем считать, что при $\ctg \alpha >\ctg \beta$, когда


$x(\gamma_{B} +0) > x(\gamma_{A}-0)$,


выполнены условия


%\begin{equation}\label{14.4}


\begin{gather*}


y\geq -b+2(x-b\ctg \beta)/ (\ctg\alpha -\ctg \beta),


\quad (x,y)\in L_{z}^{1}, \\


y\leq -b+2(x-b\ctg \beta)/ (\ctg \alpha -\ctg \beta),


\quad (x,y)\in L_{z}^{2},


\end{gather*}


%\end{equation}


при $\ctg \alpha <\ctg \beta$ --- условия


\begin{equation}\label{14.5}


\begin{split}


y&\leq -b+2(x-b\ctg \beta)/ (\ctg \alpha -\ctg \beta),


\quad (x,y)\in L_{z}^{1}, \\


y&\geq -b+2(x-b\ctg \beta)/ (\ctg \alpha -\ctg \beta),


\quad (x,y)\in L_{z}^{2},


\end{split}


\end{equation}


при $\ctg \alpha =\ctg \beta$ --- условия


\begin{equation}\label{14.6}


\begin{split}


x&\leq b\ctg \beta,\quad (x,y)\in L_{z}^{1}, \\


x&\geq b\ctg \beta,\quad (x,y)\in L_{z}^{2}.


\end{split}


\end{equation}





Остановимся на случае $\ctg \alpha =\ctg \beta$.


С учетом (\ref{14.1}), (\ref{14.2})


условие (\ref{14.6}) запишем в виде


$$


B_{0}+2A\sin\gamma\leq -2B\cos\gamma +\frac{b\ctg \beta}


{|F_{0}(e^{i\gamma})|\cos\theta\gamma}-


\frac{1}{2\pi}


\int_{0}^{2\pi}\frac{-c(\sigma)}{|F_{0}(e^{i\sigma})|}


\ctg \frac{\sigma -\gamma}{2}d\sigma,\ \; \gamma\in


(\gamma_{A},\gamma_{B}).


$$


Правая часть есть функция, ограниченная вблизи $\gamma_{A}$, $\gamma_{B}$.


Это согласуется и с тем, что данное неравенство при


$\gamma\to\gamma_{A}+0$, $\gamma\to\gamma_{B}-0$ переходит в равенство.





Последнее соотношение представим в виде


\begin{equation}\label{17.6}


B_{0}\cos\alpha_{1}(\gamma)+A\sin\alpha_{1}(\gamma)\leq E(\gamma,B),


\quad\gamma\in (\gamma_{A},\gamma_{B}),


\end{equation}


где


\begin{gather*}


%% \begin{equation}\label{17.7}


\alpha_{1}(\gamma)=\,\mathrm{arctg}\,2\sin\gamma,\\


%


E(\gamma,B)=\bigg\lbrace-2B\cos\gamma +


\frac{b\ctg \beta}{|F_{0}(e^{i\gamma})|


\cos\theta(\gamma)} -


\frac{1}{2\pi}


\int_{0}^{2\pi}\frac{-c(\sigma)}{|F_{0}(e^{i\sigma})|}


\ctg\frac{\sigma -\gamma}{2}d\sigma\bigg\rbrace/\sqrt{1+


4\sin^{2}\gamma}.


\end{gather*}





В плоскости комплексного переменного $u+iv$ неравенство


$$


u\cos\alpha_{1}(\gamma)+v\sin\alpha_{1}(\gamma)\leq E(\gamma,B),


\quad\gamma\in(\gamma_{A},\gamma_{B}),


$$


определяет семейство полуплоскостей, имеющих общую часть $D$,


поскольку для нормального вектора $(\cos\alpha_{1}(\gamma),


\sin\alpha_{1}(\gamma))$ границы каждой из полуплоскостей имеем


$$


\sup\limits_{\gamma\in(\gamma_{A},\gamma_{B})}\alpha_{1}(\gamma)-


\inf\limits_{\gamma\in(\gamma_{A},\gamma_{B})}\alpha_{1}(\gamma)=


2\sup\limits_{\gamma\in(\gamma_{A},\gamma_{B})}\alpha_{1}(\gamma)


<\pi,


$$


причем в область $D$


входит угол раствора $\pi -2\sup\limits_{\gamma\in(\gamma_{A},


\gamma_{B})}\alpha_{1}(\gamma)$ с вершиной в точке $u+iv=


\inf\limits_{\gamma\in(\gamma_{A},\gamma_{B})}E(\gamma,B)=m(B)$.


При $m(B)\geq 0$ этот угол содержит точки действительной


оси, для которых $u< m(B)$, стороны его симметричны относительно


действительной оси; при $m(B) > 0$ вершинa угла совпадает с точкой


$$


u+iv=\frac{m(B)}{\cos\big(\sup\limits_{\gamma\in(\gamma_{A},\gamma_{B})}


\alpha_{1}(\gamma)\big)}.


$$





Неравенство (\ref{17.6}), равносильное (\ref{14.6}), будет


выполняться, если постоянные $B_{0}$, $A$ будут выбраны так, что


точка $(B_{0},A)\in D$.





Считая $-\gamma_{A} <\gamma <\gamma_{A}$, совершенно


аналогично покажем, что неравенства


(\ref{14.6}) выполняются, когда постоянные $B_{0}$, $A$ выбраны так,


что $(B_{0},A)\in D_{1}$ --- области в плоскости $u+iv$, причем в


область $D_{1}$ входит угол раствора


$\pi -2\sup\limits_{\gamma\in(\gamma_{A},\gamma_{B})}\alpha_{1}(\gamma)$,


симметричный относительно


действительной оси и содержащий точки положительной полуоси, включая


сколь угодно большие значения $u$, а вершина угла совпадает


с точкой


$$


%\begin{equation}\label{19.1}


u+iv=\sup\limits_{\gamma\in(-\gamma_{A},\gamma_{A})}E(\gamma,B)=M(B)


\ \; \text{при $M(B) < 0$ }


%\end{equation}


$$


и с точкой


$$


u+iv=M(B)/\cos\big(\sup\limits_{\gamma\in(-\gamma_{A},\gamma_{A})}


\alpha_{1}(\gamma)\big) \ \; \text{при $M(B) > 0$}.


$$





Области $D$ и $D_{1}$ имеют пересечение $D_{*}$, если


\begin{alignat*}{2}


%\begin{equation}\label{19.3}


&M(B) <\frac{m(B)}{\cos(\sup\limits_{\gamma\in(\gamma_{A},


\gamma_{B})}\alpha_{1}(\gamma))}&\quad &\text{при}\ \; m(B) < 0, \\


%


% \label{19.4}


&\frac{M(B)}{\cos(\sup\limits_{\gamma\in(-\gamma_{A},


\gamma_{A})}\alpha_{1}(\gamma))} < m(B)&\quad &\text{при}\ \; m(B) > 0.


\end{alignat*}





Введем следующие обозначения:


\begin{align*}


M_{1}&=\sup\limits_{\gamma\in(-\gamma_{A},\gamma_{A})}\bigg\lbrace\bigg[


\frac{b\ctg \beta}{|F_{0}(e^{i\gamma})|}-\frac{1}{2\pi}


\int_{0}^{2\pi}\frac{-c(\sigma)}{|F_{0}(e^{i\sigma})|}


\ctg \frac{\sigma -\gamma}{2}d\sigma\bigg]/\sqrt{1+


4\sin^{2}\gamma}\bigg\rbrace, \\


%


m_{1}&=\inf\limits_{\gamma\in(\gamma_{A},\gamma_{B})}\bigg\lbrace\bigg[


\frac{b\ctg \beta}{|F_{0}(e^{i\gamma})|\cos\theta(\gamma)}-\frac{1}{2\pi}


\int_{0}^{2\pi}\frac{-c(\sigma)}{|F_{0}(e^{i\sigma})|}


\ctg \frac{\sigma -\gamma}{2}d\sigma\bigg]/\sqrt{1+


4\sin^{2}\gamma}\bigg\rbrace, \\


%


N&=\max\left\lbrace


\frac{M_{1}\cos(\sup\limits_{\gamma\in(\gamma_{A},\gamma_{B})}


\alpha_{1}(\gamma))-m_{1}}{2(1+\cos(\sup\limits_{\gamma\in(\gamma_{A},


\gamma_{B})}\alpha_{1}(\gamma)))},\


\frac{M_{1}-m_{1}\cos(\sup\limits_{\gamma\in(-\gamma_{A},\gamma_{A})}


\alpha_{1}(\gamma))}{2(1+\cos(\sup\limits_{\gamma\in(-\gamma_{A},


\gamma_{A})}\alpha_{1}(\gamma)))}\right\rbrace.


\end{align*}


Теперь выберем постоянную $B$ так, чтобы $0 > B > N$, тогда


области


$D$ и $D_{1}$ имеют непустое, зависящее от $B$, пересечение $D_{*}$;


выберем числа $B_{0}$, $A$ так, чтобы $(B_{0},A)\in D_{*}$,


тогда будут выполняться неравенства (\ref{14.6}). Если


при выбранных значениях постоянных $B_{0}$, $A$, $B$ имеет место


неравенство (\ref{14.3}), то область $D_{z}$ будет однолистной.





Неравенство (\ref{14.5}) запишем так


(для $\ctg \alpha <\ctg \beta$):


$$


B_{0}\cos\alpha_{1}(\gamma)+A\sin\alpha_{1}(\gamma)\geq\wt{E}


(\gamma,B),\quad -\gamma_{A} <\gamma <\gamma_{A},


$$


где $\alpha_{1}(\gamma)$ обозначает то же, что и выше,


\begin{multline*}


\wt{E}(\gamma,B)=\bigg\lbrace -2B\cos\gamma +\bigg[b\ctg \beta +


\frac{(\ctg \alpha -\ctg \beta)(y(\gamma)+b)}{2}\bigg]


/|F_{0}(e^{i\gamma})|- \\


%


-\frac{1}{2\pi}


\int_{0}^{2\pi}\frac{-c(\sigma)}{|F_{0}(e^{i\sigma})|}


\ctg \frac{\sigma -\gamma}{2}d\sigma\bigg\rbrace/\sqrt{1+


4\sin^{2}\gamma}.


\end{multline*}


Это неравенство выполняется,


когда постоянные $B_{0}$, $A$ выбраны так, чтобы точка $(B_{0},A)\in


\wt{D}$ --- области в плоскости $u+iv$,


причем в область $\wt{D}$ входит угол раствора $\pi -


2\sup\limits_{\gamma\in(-\gamma_{A},\gamma_{A})}\alpha_{1}(\gamma)$,


симметричный относительно действительной оси и содержащий точки


положительной полуоси,


включая сколько угодно большие значения $u$, а вершина угла


совпадает с точкой


$$


u+iv=\sup\limits_{\gamma\in(-\gamma_{A},\gamma_{A})}\wt{E}(\gamma,


B)=\wt{M}(B)\quad \text{при}\ \; \wt{M}(B) < 0,


$$


с точкой


$$


u+iv=\wt{M}(B)/\cos(\sup\limits_{\gamma\in(-\gamma_{A},\gamma_{A})}


\alpha_{1}(\gamma))\quad \text{при}\ \;\wt{M}(B) > 0.


$$


Аналогично убедимся в том, что неравенство (\ref{14.3}) выполняется,


когда


постоянные $B_{0}$, $A$ выбраны так, чтобы точка $(B_{0},A)\in\wh{D}$


--- oбласти в плоскости $u+iv$, причем в область $\wh D$


входит угол раствора $\pi -2\sup\limits_{\gamma\in(\gamma_{A},


\gamma_{B})}\alpha_{1}(\gamma)$,


симметричный относительно действительной оси и содержащий точки


положительной действительной оси, а вершина угла совпадает с точкой


$$


u+iv=\sup\limits_{\gamma\in(-\gamma_{A},\gamma_{A})}\bigg\lbrace


-B\cos\gamma-\frac{1}{2\pi}


\int_{0}^{2\pi}\frac{-c(\sigma)}{|F_{0}(e^{i\sigma})|}


\ctg \frac{\sigma -\gamma}{2}d\sigma\bigg\rbrace/\sqrt{1+


4\sin^{2}\gamma}=


\wh{M}(B)\ \; \text{при}\ \;\wh{M}(B) < 0,


$$


с точкой


$$


u+iv=\wh{M}(B)/\cos(\sup\limits_{\gamma\in(\gamma_{A},\gamma_{B})}


\alpha_{1}(\gamma))\quad \text{при}\ \;\wh{M}(B) > 0.


$$





Ясно, что области $\wt D$, $\wh D$ имеют непустое пересечение


$\wt D_{*}$, содержащее, в частности, пересечение


вышеуказанных углов, входящих в $\wt D$, $\wh D$.


Пусть при заданном $B$ постоянные $B_{0}$, $A$


выбраны так, что $(B_{0},A)\in\wt D_{*}$, тогда имеют место


неравенства (\ref{14.3}), (\ref{14.5})


и область $D_{z}$ будет однолистной.





Аналогичным образом исследуется случай


$\ctg \alpha>\ctg \beta$. И здесь 


при некотором фиксированном $B$


имеется непустая область $\wh D_*$ такая,


что при $(B_0,A)\in \wh D_*$ искомое 


решение является однолистным.


Выводом из этих рассуждений является 





\begin{thm}


Решение $z(\zeta)$ ОКЗ по параметрам $\theta$, $y$, 


$-\beta\le\theta\le\alpha$, определяется по формуле $(\ref{4.4})$.


Это решение будет однолистным, если входящие в $(\ref{4.4})$


произвольные постоянные удовлетворяют неравенству $(\ref{14.3})$


и точка $(B_0,A)\in D_*$ при $\ctg \alpha=\ctg \beta$ 


либо $(B_0,A)\in \wt {D_*}$ при $\ctg \alpha<\ctg \beta$,


либо $(B_0,A)\in \wh {D_*}$ 


при $\ctg \alpha>\ctg \beta$.


\end{thm}








\begin{thebibliography}{9}





\bibitem{TN}


Тумашев Г.Г., Нужин М.Т. {\em Обратные краевые задачи и их


приложения}. -- 2-е изд. -- 


Казань: Изд-во Казанск. ун-та, 1965. -- 333 с.


\bibitem{Mu}


Мусхелишвили Н.И. {\em Сингулярные интегральные уравнения. Граничные


задачи теории функций и некоторые их прилoжения к математической


физике}. -- М.: Наука, 1968. -- 511 с.


\bibitem{AS}


Абубакиров Н.Р., Салимов Р.Б. {\em Новый подход к решению краевой задачи 


Гильберта для аналитической функции в многосвязной области}. --


Казанск. ун-т. -- Казань, 1999. -- 


15 с. -- Деп. в ВИНИТИ 28.05.99, \No\,1703-B99. 


\bibitem{SSh}


Салимов Р.Б., Шабалин П.Л. {\em Новый подход к решению краевой задачи


Гильберта с разрывными коэффициентами для аналитической в круге


функции} // Наука и язык. -- Казань, 1999. -- С.\,33--37.


\bibitem{Ga}


Гахов Ф.Д. {\em Краевые задачи}. -- М.: Наука, 1977. -- 640 с.


\bibitem{Sa}


Салимов Р.Б. {\em К вычислению сингулярных интегралов с ядром Гильберта} //


Изв. вузов. Математика. -- 1970. -- \No\,12. -- C.\,93--96.





\end{thebibliography}





\vskip 0.5cm





\post{\it Казанский государственный университет}{\it Поступила}


\post{\it Казанская государственная}{07.07.2000}


\post{\it архитектурно-строительная академия}{}





\label{end}


\end{document}


